\section{Stateful Agents}
The models discussed thus far are \textbf{stateless agents}. 
This means they do not maintain a memory of previous interactions; instead, they treat every request as an independent event and respond based solely on the current input.

To illustrate this, imagine you have access to an infinite sequence of \textit{superhumans}. 
You approach each one and ask exactly one question. While each superhuman provides a brilliant response 
to that specific query, they possess no knowledge of your previous conversations with their predecessors. 
Every time you move to a new person, the context is reset to a blank slate. This behavior characterizes a stateless system.

In contrast, a \textbf{stateful agent} retains information across multiple turns. 
It stores previous inputs and outputs—often referred to as the \textit{conversation history} or \textit{context}—allowing 
it to understand follow-up questions and maintain continuity in a dialogue.

In this section we will explore various techniques used in agents design, from simple stateless agents to more complex stateful ones.

\subsection{Stateless}
As mentioned, stateless agents do not maintain any memory of past interactions. The \texttt{chat completition}
api is an example of a stateless agent by design.
\begin{lstlisting}[language=Python]
# Turn 1: tell the model your name
r1 = client.chat.completions.create(
    model="gpt-4.1-mini",
    messages=[{"role": "user", "content": "My name is Alice."}],
)
# Turn 1: "Nice to meet you, Alice!"

# Turn 2: ask for your name  but we DON'T send turn 1
r2 = client.chat.completions.create(
    model="gpt-4.1-mini",
    messages=[{"role": "user", "content": "What is my name?"}],
)
# Turn 2: "I don't have access to your name."
\end{lstlisting}
The model can't answer because it literally never saw the first message.

Although this type of agents seems pretty useless, they can be fine for single-shot tasks,
like writing a poem, generating an image, or answering a question.

\subsection{Default Stateful Agents}
Open AI \texttt{Response API} can be used to create stateful agents.
\begin{lstlisting}[language=Python]
resp1 = client.responses.create(model="gpt-4.1-mini", input="My name is Alice.")

resp2 = client.responses.create(
    model="gpt-4.1-mini",
    input="What is my name?",
    previous_response_id=resp1.id,   # server remembers turn 1
)
# "Your name is Alice."
\end{lstlisting}
This is convenient, but it's a black box approach. We have no control over what the model
remembers and how it uses that information.

All the next techniques we will manage the state ourselves, giving us more control and flexibility.

\subsection{Full History}
A first approach to create a stateful agent is to send with each request the full conversation history.
This is obviously the most straightforward way to maintain context, but it can quickly become inefficient as the conversation grows longer.
In fact, each time we send a new message, the number of tokens we need to send increases quadratically, since we need to resend all previous messages.

This is not only inefficient in terms of latency and cost, but it can also lead to fulling the model's context window,
which can cause lower quality responses or even errors.

\subsection{Sliding Window}
To mitigate the issues of the full history approach, we can use a \textbf{sliding window} technique.
Instead of sending the entire conversation history, we only send the most recent $N$ messages, where $N$ is a predefined window size.
This way, we can maintain some context while keeping the number of tokens manageable and constant over time.

The downside of this approach is that we might lose important information from earlier in the conversation, 
which can lead to less coherent responses if the model relies on that information.

\subsection{Summarization}
To address the limitations of the sliding window technique, we have to find a way to retain important information from the conversation without sending the entire history.
This is where \textbf{summarization} comes into play.
The idea is to periodically summarize the conversation history and send the summary along with the most recent messages. 
This way, we can keep the context manageable while still retaining important information from earlier in the conversation.

The summarization is usually done by the model itself. Along with the user message,
we can provide the model with the current summary and ask it to update the summary based on the new message.

This way, the memory stays bounded, without the loss of information.

A few tips for effective summarization prompts:
\begin{itemize}
    \item Use a structured format for the summary, organizing information into categories (e.g., "Facts", "Opinions", "Actions", etc.);
    \item Set a maximum length for the summary to ensure it remains concise and focused on the most relevant information;
    \item Resolve contradictions in the summary by asking the model to identify and address any conflicting information.
\end{itemize}

\subsection{Long-Term Memory}
Every technique we have seen so far is based on the idea of a \textit{working memory}, that dies whenever the conversation ends.
However, in many applications we want to have a \textit{long-term memory} that persists across multiple conversations and sessions. 

This is usually achieved by:
\begin{itemize}
    \item Once the conversation ends, extract important information. Ask the LLM to pull out structured facts worth remembering: facts, preferences, decisions, expertise, etc;
    \item Store this information in an external database;
    \item When a new conversation starts, retrieve relevant information from the database and inject it into the prompt.
\end{itemize}
This is the actual pattern used behind the OpenAI \texttt{ChatGPT} memory feature.

\subsection{When to use which pattern}
Each strategy can be used in different scenarios, depending on the specific requirements of the application and the constraints of the model being used.
\begin{itemize}
    \item \textbf{Stateless}: suitable for single shot tasks (e.g., translate, classify);
    \item \textbf{Full History}: suitable for short conversations;
    \item \textbf{Sliding Window}: suitable for longer conversations where only the most recent context is relevant;
    \item \textbf{Summarization}: suitable for long conversations where important information needs to be retained without exceeding the context window;
    \item \textbf{Long-Term Memory}: suitable for applications that require personalization or have returning users, where information needs to persist across sessions.
\end{itemize}